\section{Estudos sobre \gls{rvdb} e aproximação relativa}
\label{sec_evolucao_rendezvous_aproximacao_relativa}

Historicamente, o estudo do movimento relativo entre espaçonaves e o planejamento de manobras de aproximação consolidaram-se a partir da formulação clássica de \cite{clohessy1960terminal}. 
No entanto, ao longo do desenvolvimento da área, observou-se que o problema de \gls{rvdb} não poderia permanecer restrito apenas à dinâmica relativa linearizada.
Em missões de serviço em órbita, inspeção, captura e acoplamento, passou a ser necessário considerar não apenas o movimento relativo geométrico entre perseguidor e alvo, mas também restrições operacionais, requisitos de segurança e diferentes níveis de cooperação do alvo durante a operação.

Nesse contexto, \cite{arantes2010far} apresentam um estudo sobre manobras \textit{far} e \textit{proximity} de uma constelação de satélites de serviço, tratando de forma integrada diferentes fases da missão. Os autores empregam a dinâmica relativa baseada nas equações de \gls{hcw} e discutem tanto as manobras distantes quanto as operações de proximidade. O trabalho incorpora ainda a estimação autônoma da pose do satélite cliente, mostrando que, em missões de serviço orbital, a aproximação não depende apenas da formulação da trajetória relativa, mas também da capacidade de acompanhar a configuração espacial do alvo. Os resultados apresentados pelos autores mostram a viabilidade da abordagem proposta para integrar análise orbital e operação de proximidade em um mesmo contexto de \gls{rvdb}.

Mais recentemente, \cite{fiori2024modeling} apresentam um arcabouço de modelagem, simulação e controle para \gls{rvdb} automatizado nas proximidades de uma estação espacial, considerando restrições posicionais e obstáculos físicos. Nesse trabalho, o movimento da espaçonave é formulado por meio de um formalismo geométrico baseado em grupos de Lie, enquanto o problema de aproximação é tratado com o uso de potenciais virtuais. Os autores realizam simulações numéricas com o objetivo de comparar estratégias de controle e avaliar seu desempenho no cenário estudado.

De modo geral, a literatura revisada mostra que os estudos sobre \gls{rvdb} e aproximação relativa evoluíram de formulações clássicas centradas na dinâmica relativa linearizada para abordagens mais amplas, nas quais passam a ser considerados aspectos operacionais, geométricos e autônomos da missão. Enquanto os trabalhos clássicos estabeleceram a base matemática para o tratamento do movimento relativo, estudos posteriores ampliaram o escopo do problema ao incorporar requisitos de segurança, interação com alvos cooperativos e não cooperativos, além de integração com tarefas de navegação e operação de proximidade. Esses aspectos são relevantes para a presente dissertação, pois mostram que a fase de aproximação curta deve ser analisada dentro de um contexto mais amplo, no qual dinâmica relativa, atitude, restrições operacionais e interação com o alvo tornam-se progressivamente acopladas.